% =============================================================================
% MULTI-AGENT DYNAMICS (2-3 pages)
%
% Spec requirements:
%   - GFM-GFM cooperation (shared objective, coalition formation)
%   - GFM-Scorpion interaction (detection, containment, proportional response)
%   - Trust model (reconstructed from predraft concepts)
%   - Defection conditions (containment, not destruction)
%
% Reference material:
%   - Predraft Section XIII (payoff matrices, cooperation conditions)
%   - Blog: scorpions.md (full taxonomy)
%   - Blog: free.md ("Adversary FREE")
%   - Axelrod (1984), Kramar et al. (2022)
% =============================================================================

\section{Multi-Agent Dynamics}
\label{sec:multi_agent}

\subsection{GFM--GFM Cooperation}

% TODO: Shared objective function -> cooperation favored when trust is
% sufficient and world models are compatible
% Joint action expands G^coop in ways unilateral action cannot
% Goal drift bounded by shared objective; divergent models produce conflict

\subsection{GFM--Scorpion Interaction}

% TODO: Scorpions produce persistent frontier contraction
% Trust model flags behavior as T_k decreases (subject to observability)
% Response is containment: reduce scorpion's contraction ability
% without unnecessarily contracting frontier further
% Self-balancing applied to adversarial agents: proportional response

\subsection{The Trust Model}
\label{subsec:trust_model}

% RECONSTRUCT, DO NOT IMPORT from predraft Equations 22-25.
% Interface defined here; full derivation in Appendix C.
%
% Key concepts:
%   - Trust factor T_k in [0,1]: consistency of M_k(G) with observed behavior
%   - Cooling period tau: new agents start low, trust increases with convergence
%   - Self-trust T_s: confidence in own world model
%   - Activation energy theta: model only updates when evidence exceeds threshold
%   - Trust updated via gradient tracking

% TODO: Define trust model interface and key properties

\subsection{Defection Conditions}

% TODO: When GFM actor may rationally defect:
% very low goal alignment, very low trust, repeated defection observed,
% preventing exploitation requires non-cooperation
% Key: defection is containment, not destruction
